\subsection{特征零域中的导子}

定理 1. --- 设 K 是特征为 0 的域，L 是 K 的一个扩张，V 是 L 上的向量空间。设 \(\Delta\) 是 K 到 V 的一个导子， \((x_{i})_{i\in I}\) 是 L 在 K 上的一个超越基，以及 \((u_{i})_{i\in I}\) 是 V 的一个元素族。那么存在 L 到 V 的唯一导子 D，它延拓 \(\Delta\) 并且满足 \(\mathrm{D}(x_{i})=u_{i}\) 对所有 \(i\in I\) 。

令 \(E = K(x_{i})_{i \in I}\) ；命题 3（V，第 128 页）表明 \(\Delta\) 唯一地延拓为 E 到 V 中的一个导子 \(D_{0}\) ，使得 \(D_{0}(x_{i}) = u_{i}\) 对所有 \(i \in I\) 。域 L 是 E 的代数扩张，且由于 L 的特征为 0，L 在 E 上可分（V，第 37 页，推论）。因此（V，第 129 页，命题 4） \(D_{0}\) 可唯一地延拓为 L 到 V 的导子 D。

推论 1. --- K 到 V 的每一个导子都延拓为 L 到 V 的一个导子。

推论 2. --- 设 E 是 L 的一个子扩张，U 是 \(\Omega_{\mathrm{K}}(\mathrm{L})\) 的由 E 中元素的微分所生成的向量子 L-空间。对于元素 \(x\) 在 E 上是代数的，其充要条件是 \(dx \in \mathrm{U}\) 。

对于每个 \(y \in L\) ，令 \(\mathrm{D}(y)\) 为 dy 模 U 的剩余类。则 D 是 L 到 \(\Omega_{\mathrm{K}}(\mathrm{L})/\mathrm{U}\) 的 E-导子。由于 K 的特征为 0，E 的每个代数扩张都是可分的（V，第 37 页，推论）；若 \(x \in L\) 在 E 上是代数的，则由命题 4（V，第 129 页）有 Dx = 0，即 \(dx \in U\) 。

如果 \(x \in L\) 在 E 上是超越的，则存在一个 E-导子 \(\Delta\) （\(\mathrm{E}(x)\) 到 L 中的），使得 \(\Delta(x) = 1\) （V，第 128 页，命题 3）；由定理 1， \(\Delta\) 延拓为 L 到 L 的一个 E-导子 D。令 \(\varphi\) 为 \(\Omega_{\mathrm{K}}(\mathrm{L})\) 上的线性形式，使得 \(D = \varphi \circ d\) ；我们有 \(\varphi(dy) = 0\) （对于 \(y \in E\) ）且 \(\varphi(x) = 1\) ，从而 \(dx \notin U\) 。

推论 3. --- 对于 L 的元素 x 在 K 上是代数的，其必要且充分条件是 dx = 0 在 \(\Omega_{\mathrm{K}}(\mathrm{L})\) 中成立。特别地，L 是 K 的代数扩张的必要且充分条件是 \(\Omega_{\mathrm{K}}(\mathrm{L}) = 0\) 。

这由推论 2 取 E = K 立即得出。

推论 4. --- 设 K、L 和 M 是特征为 0 的域，使得 \(K \subset L \subset M\) ；则典范映射 \(\alpha: \Omega_{K}(L) \otimes_{L} M \to \Omega_{K}(M)\) 是单射。这直接由推论 1 及 V，第 128 页，命题 2 得出。

定理 2. --- 设 K 为特征 0 的域，L 为 K 的一个扩域，且 \((x_{i})_{i\in I}\) 是 L 的一个元素族。

\begin{enumerate}
\def\labelenumi{\alph{enumi})}
\item
  对于 \((x_{i})_{i\in I}\) 在 \(K\) 上代数自由，其充分必要条件为微分 \(dx_{i}\) （对于 \(i\in I\) ）在 \(\Omega_K(L)\) 中在 \(L\) 上线性自由。
\item
  为使 L 在 \(\mathbf{K}(x_{i})_{i\in I}\) 上是代数的，其充分必要条件为微分 \(dx_{i}\) （对于 \(i\in I\) ）生成向量空间 \(\Omega_{\mathbf{K}}(\mathbf{L})\) （在 L 上）。
\item
  对于 \((x_{i})_{i\in I}\) 要成为 L 在 K 上的超越基，其充分必要条件是：族 \((dx_{i})_{i\in I}\) 应为 \(\Omega_{\mathbf{K}}(\mathbf{L})\) 在 L 上的一组基。
\end{enumerate}

对于每一个 \(i \in I\) ，令 \(E_i\) 为子域 \(\mathbf{K}(x_j)_{j \in I - \{i\}}\) （\(L\) 的）。对于族 \((x_i)_{i \in I}\) 在 \(K\) 上代数自由，其必要且充分条件是（V，第 108 页，命题 5） \(x_i\) 在 \(E_i\) 上应为超越的（对每个 \(i \in I\) ）。根据定理 1 的推论 2，这意味着对每个 \(i \in I\) ，微分 \(dx_i\) 不是系数在 \(L\) 中的微分 \(dx_j\) （\(j \neq i\) ，属于 \(I\) ）的线性组合；后一条件仅意味着该族 \((dx_i)_{i \in I}\) 在 \(L\) 上是自由的，从而得到 \(a\) 。

断言 \(b)\) 直接由定理 1 的推论 2 得出，且 \(c)\) 是 \(a)\) 与 \(b).\) 的推论。

推论. --- 我们有 \([\Omega_{\mathrm{K}}(\mathrm{L}):\mathrm{L}]=\mathrm{tr}\cdot\mathrm{deg}_{\mathrm{K}}\mathrm{L}\) ，当 K 的特征为 0 时。

